\section{Introduction}

Large language model (LLM) agents are rapidly evolving from conversational
assistants into autonomous system operators capable of executing shell commands,
manipulating filesystems, launching subprocesses, and interacting with complex
software environments. This autonomy introduces a fundamental security
challenge: hallucinations, prompt injections, compromised context, or malicious
tools can be translated into concrete system actions at machine speed. We adopt
a strict threat model: the host operating system and \SystemName's trusted
runtime components are trusted, while the agent framework, its generated
instructions, permission declarations, and invoked binaries are treated as
potentially malicious.

Existing agent deployments often grant tools broad filesystem and network
access. In such settings, a single successful injection can exfiltrate sensitive
files, corrupt project state, or interact with remote services on the user's
behalf. Securing agents therefore requires authorization mechanisms that
precisely constrain resource accesses according to the intended task. Because
agents routinely invoke shells, interpreters, package managers, and arbitrary
binaries, these constraints must ultimately be enforced at the operating-system
boundary rather than only within the application framework.

The central difficulty is the semantic gap between agent intent and kernel
enforcement. Operating systems reason about files, sockets, capabilities,
processes, and system calls, whereas agent objectives are expressed through
natural-language instructions and tool invocations. Existing approaches attempt
to bridge this gap using task-level sandboxing, LLM-generated policies, or
per-invocation validation, but they face a three-way tension among semantic
precision, enforcement strength, and execution latency. Coarse task-level
policies fail to capture dynamic runtime behavior; per-invocation semantic
validation improves precision but places LLM reasoning on the critical path; and
kernel isolation provides strong enforcement but lacks task-level context.

This paper argues that these properties need not be tightly coupled. Our key
insight is to separate \emph{kernel-enforced confinement} from \emph{semantic
release authorization}. Before untrusted tool code runs, the system installs a
deterministic kernel policy that bounds the tool's authority. The tool then
executes speculatively within that confined environment while a trusted semantic
validator independently decides whether its outputs and effects may be released.
Reversible local effects are isolated until approval, while irreversible effects,
such as network transmissions and API calls, are treated as release boundaries
and withheld until authorization succeeds. This hides semantic validation latency
without creating an unconstrained execution window.

We present \SystemName, a runtime security framework that combines semantic
authorization, kernel-level enforcement, and speculative execution.
\SystemName treats agent-supplied permission declarations as untrusted input. A
deterministic compiler translates each declaration into a bounded, auditable,
kernel-enforceable policy, which the runtime installs before transferring
control to untrusted tool code. The tool executes under eBPF/LSM confinement,
while a trusted semantic validator concurrently evaluates the task context,
accepted permissions, and pending release request. Outputs, filesystem changes,
and externally visible effects are released only after semantic approval;
otherwise, speculative state is discarded and the invocation-scoped policy is
revoked.

\SystemName consists of three core mechanisms. First,
\emph{Semantic-to-Kernel Intermediate Representation (SK-IR)} deterministically
translates untrusted permission declarations into kernel-enforceable upper
bounds on runtime authority. SK-IR does not decide whether a permission is
semantically justified; instead, it validates, normalizes, and narrows requested
authority, rejecting declarations that cannot be represented safely.

Second, \emph{ghost policies} provide transient invocation-scoped enforcement
contexts. \SystemName materializes SK-IR outputs as eBPF/LSM policies indexed by
fresh cgroup v2 contexts, ensuring that tools and their descendants remain
subject to the same least-privilege policy. Ghost policies are lease-bound and
revoked when the invocation completes, fails, rolls back, times out, or expires.

Third, \emph{speculative execution with ShadowFS} overlaps confined execution
with semantic validation. Filesystem mutations are redirected into
per-invocation copy-on-write overlays, and agent-visible outputs are buffered
until approval. If authorization succeeds, approved state is committed and
outputs are released; otherwise, overlays and buffered outputs are discarded
before becoming host-visible. Irreversible operations are routed through
mediated release paths and withheld until semantic approval succeeds.

\FIXME{We implement \SystemName in Python, Go, and C on Linux. Evaluated against
agent-security attacks including prompt injections, malicious tools, and
privilege-escalation attempts, \SystemName blocks unauthorized actions with only
\textbf{[X]}\% throughput overhead, while hiding up to \textbf{[Y]}\% of
authorization latency (avg. \textbf{[Z]}s) for filesystem-intensive workloads.}

In summary, this paper makes the following contributions:
\begin{itemize}
\item SK-IR, a deterministic semantic-to-kernel compilation framework that
translates untrusted agent permission declarations into auditable,
kernel-enforceable bounds on runtime authority.

\item Ghost policies, transient invocation-scoped eBPF/LSM security contexts
that bind least-privilege enforcement to individual tool executions and revoke
authority automatically.

\item Speculative execution with ShadowFS, which decouples semantic validation
from tool execution by isolating reversible effects, buffering outputs, and
mediating irreversible release boundaries.

\item An implementation and evaluation of \SystemName, demonstrating that
semantic release authorization can be combined with kernel-grade confinement
while reducing authorization latency on the critical path.
\end{itemize}